% Owner: Duong. Target for Abstract, Introduction, Related Work, and Table I: 1,700--1,900 words at completion.
\section{Introduction}\label{sec:introduction}

\subsection{Background and Motivation}
Rental property management is a chain of related decisions rather than a single record-keeping task. A listing refers to a physical property; a rental request may precede a contract; an active tenancy then generates billing, utility accounting, and maintenance obligations. Each stage involves different actors and a different data lifecycle. In a centralized implementation that places these changes in one database transaction, the database can provide atomicity and referential integrity for that transaction. A shared schema couples schema changes across modules, while a request path that waits for several components is exposed to delays or failures in those components.

Splitting the workflow into independently owned services can reduce shared-schema coupling but introduces distributed state. A systematic mapping classified 103 primary microservice studies from 532 candidates \cite{dipfrancesco2019architecting}. Boundary derivation remains difficult: a review of 34 studies and a survey of 63 practitioners report varied interpretations of domain-modeling artifacts \cite{zhong2024ddd}. A mixed-method study of microservice data management identifies ownership, fault and performance isolation, and independent schema evolution as recurring concerns \cite{laigner2021datamanagement}.

In the proposed rental workflow, events propagate committed local changes, but downstream state may lag the initiating response. Mining more than 8,000 Stack Overflow questions, Laigner et al.\ identify difficulties with payload size, schema modeling, flow debugging, and ordering \cite{laigner2026eventmanagement}. In one application comparison, a monolith used fewer resources and had better response times than an event-driven implementation \cite{cabane2024edaperformance}. Neither result predicts this platform's performance. The present work therefore specifies what would need to be measured before making such a claim.

\subsection{Problem Statement}
The central problem is how to coordinate a rental workflow when participating services own separate records and may fail independently. A successful acceptance response would state only that Listing and Booking committed its local transition; it would not establish contract creation, notification, or later billing. A redelivered acceptance event could create another draft unless Contract recognizes the prior effect. With separate local transactions, intermediate states can be visible before the workflow terminates. Saga literature describes compensation as a semantic undo rather than restoration of a prior database snapshot \cite{garciamolina1987sagas,stefanko2019saga}; this paper does not claim to have implemented a complete compensating saga.

Two constraints follow. Authenticating a caller at the entry point does not authorize an operation on a resource another service owns; the service making the resource decision needs an ownership or building-assignment check. The transactional outbox has been described for an applied redesign and as a Spring implementation methodology, without controlled reliability results, while service-level idempotence does not by itself establish broker-consumer deduplication \cite{peliza2025outbox,glybovets2024transactions,ramalingam2013idempotence}. This paper therefore specifies tests for these mechanisms in the rental workflow.

\subsection{Research Objectives and Questions}
Three objectives follow: define the proposed business boundaries and data ownership for seven services; specify how synchronous reads and asynchronous events would cooperate without implying a distributed atomic transaction; and define reproducible tests for ownership, latency, duplication, and recovery under controlled faults.

These motivate three questions. RQ1: which service owns each rental decision and record, and what source and database evidence would show that the proposed boundaries avoid cross-service table writes? RQ2: under specified retry and failure points, does each committed acceptance yield at least one corresponding publication and exactly one correct draft by a predefined cutoff? RQ3: what API latency, throughput, and recovery times does the implementation exhibit under a stated workload? Section~\ref{sec:architecture} gives a provisional design answer to RQ1; Section~\ref{sec:evaluation} defines checks for all three questions. RQ2 and RQ3 remain unanswered empirically.

\subsection{Main Contributions}
The paper contributes a proposed service-ownership model and identifies the records that cross boundaries by identifier; a candidate acceptance-to-contract event workflow with its local commit and deduplication points; and a protocol that classifies completed, pending, dead-lettered, and unexplained outcomes under load and injected failures. A source audit identifies gaps between the target topology and the supplied backend snapshot. These artifacts support design review and future measurement, not claims of comparative performance or production reliability.

\subsection{Paper Organization}
Section~\ref{sec:related} states the gap this study addresses. Section~\ref{sec:architecture} presents service boundaries, data ownership, and access control. Section~\ref{sec:events} specifies communication and consistency mechanisms. Section~\ref{sec:evaluation} gives the measurement protocol; Sections~\ref{sec:discussion} and~\ref{sec:conclusion} discuss and close the paper.
